%% toc: P3 24 24.3
\section{Proofs for definable sets with parameters}
\label{sec:annex-p03-w03-p03-003-definable-sets}

Notation is that of \cref{def:p03-w03-p03-003-declared-formula} and
\cref{def:p03-w03-p03-003-coordinate-maps}; \(\mathcal K\) ranges over
\(\Lang\)-structures with domain \(K\), and \(s\) is an arbitrary assignment
in the structure at hand.  For a declared formula \(\varphi(\bar x;\bar y)\),
\(\bar a\in K^n\) and \(\bar b\in K^m\), the definition and
\cref{lem:fol-coincidence} give
\[
 \bar a\in\varphi(K^n,\bar b)
 \iff
 \mathcal K,s[\bar x\mapsto\bar a,\bar y\mapsto\bar b]\entails\varphi
 \quad\text{for every assignment }s.
 \tag{\(\ast\)}
\]
Three facts are used throughout.

\emph{(F1) Blocks of quantifiers.}  Let \(\bar v=(v_1,\ldots,v_p)\) be
pairwise distinct variables, \(p\geq0\).  Then
\(\mathcal K,t\entails\exists v_1\cdots\exists v_p\,\theta\) iff
\(\mathcal K,t[\bar v\mapsto\bar e]\entails\theta\) for some
\(\bar e\in K^p\), and \(\mathcal K,t\entails\forall v_1\cdots\forall
v_p\,\theta\) iff this holds for every \(\bar e\in K^p\).  For \(p=0\) there
is nothing to prove.  For \(p>0\), the satisfaction clause for
\(\exists v_1\) reduces the left side to: for some \(e_1\),
\(\mathcal K,t[v_1\mapsto e_1]\entails\exists v_2\cdots\exists v_p\,\theta\);
by induction on \(p\) this says that for some \(e_1,\ldots,e_p\),
\(\mathcal K,t[v_1\mapsto e_1][v_2\mapsto e_2]\cdots[v_p\mapsto e_p]
\entails\theta\), and the iterated update is \(t[\bar v\mapsto\bar e]\)
because the \(v_i\) are distinct.  The universal case is the same with
``every''.

\emph{(F2) Reducts.}  If \(\Lang\subseteq\Lang'\), \(\mathcal K'\) is an
\(\Lang'\)-structure and \(\varphi\) an \(\Lang\)-formula, then
\(\mathcal K',t\entails\varphi\) iff
\(\mathcal K'\mathbin{\upharpoonright}\Lang,t\entails\varphi\) for every
assignment \(t\).  Indeed both structures have the same domain and the same
interpretation of every symbol of \(\Lang\), so every \(\Lang\)-term has the
same value under \(t\) in both, by induction on terms; hence every atomic
\(\Lang\)-formula has the same truth value; the connective clauses preserve
this, and the quantifier clauses range over the same domain.

\emph{(F3) Equations pin values.}  If \(t\) is an assignment and \(u,v\) are
variables, then \(\mathcal K,t\entails u=v\) iff \(t(u)=t(v)\), by the
satisfaction clause for equality and the value of a variable.

\begin{annexitem}{annex:p03-w03-p03-003-normal-form}{thm:p03-w03-p03-003-normal-form}{Renaming, padding and constants for parameters}
\begin{proof}
\emph{Renaming.}  Write \(\bar v=(v_1,\ldots,v_p)=(\bar x,\bar y)\) and
\(\bar v'=(v'_1,\ldots,v'_p)=(\bar u,\bar w)\), with \(p=n+m\geq1\).  The
variable set is infinite (\cref{sec:fol-syntax-semantics-01}), so we may
choose pairwise distinct variables \(z_1,\ldots,z_p\) outside the finite set
\(\operatorname{Var}(\varphi)\cup\{\bar v\}\cup\{\bar v'\}\).  Put
\[
 \varphi_1\defeq\exists v_1\cdots\exists v_p
   \Bigl(\bigwedge_{i\leq p}v_i=z_i\land\varphi\Bigr),\qquad
 \varphi'\defeq\exists z_1\cdots\exists z_p
   \Bigl(\bigwedge_{i\leq p}z_i=v'_i\land\varphi_1\Bigr).
\]
These depend only on \(\varphi\) and the two lists (the choice of the
\(z_i\) is syntactic).  By the free-variable clauses,
\(\FV(\varphi_1)\subseteq(\{\bar v\}\cup\{\bar z\}\cup\FV(\varphi))
\setminus\{\bar v\}\subseteq\{\bar z\}\), since
\(\FV(\varphi)\subseteq\{\bar v\}\) and no \(z_i\) is a \(v_j\).  Likewise
\(\FV(\varphi')\subseteq(\{\bar z\}\cup\{\bar v'\})\setminus\{\bar z\}
=\{\bar v'\}\), since no \(z_i\) is a \(v'_j\).

Let \(t\) be any assignment in \(\mathcal K\) and write
\(t(\bar z)=(t(z_1),\ldots,t(z_p))\), similarly for other lists.  By (F1),
\(\mathcal K,t\entails\varphi_1\) iff
\(\mathcal K,t[\bar v\mapsto\bar e]\entails\bigwedge_iv_i=z_i\land\varphi\)
for some \(\bar e\in K^p\).  No \(z_i\) is updated, so by (F3) the equations
hold iff \(\bar e=t(\bar z)\).  Thus
\[
 \mathcal K,t\entails\varphi_1\iff
 \mathcal K,t[\bar v\mapsto t(\bar z)]\entails\varphi .
\]
The same argument, with \(\bar z\) quantified and no \(v'_j\) updated, gives
\(\mathcal K,t\entails\varphi'\) iff
\(\mathcal K,t[\bar z\mapsto t(\bar v')]\entails\varphi_1\).  Applying the
first equivalence to \(t'=t[\bar z\mapsto t(\bar v')]\), whose values on
\(\bar z\) are \(t(\bar v')\), we get
\(\mathcal K,t\entails\varphi'\) iff
\(\mathcal K,t'[\bar v\mapsto t(\bar v')]\entails\varphi\).  The last
assignment sends \(\bar v\) to \(t(\bar v')\), so by
\cref{lem:fol-coincidence} this holds iff
\(\mathcal K,s[\bar v\mapsto t(\bar v')]\entails\varphi\) for any \(s\).
Now take \(t=s[\bar u\mapsto\bar a,\bar w\mapsto\bar b]\), so that
\(t(\bar v')=(\bar a,\bar b)\).  By \((\ast)\) for both declared formulas,
\(\bar a\in\varphi'(K^n,\bar b)\) iff \(\bar a\in\varphi(K^n,\bar b)\).

\emph{Padding.}  The list \((\bar x;\bar y\bar z)\) consists of distinct
variables and contains \(\FV(\varphi)\).  The assignments
\(s[\bar x\mapsto\bar a,\bar y\mapsto\bar b,\bar z\mapsto\bar c]\) and
\(s[\bar x\mapsto\bar a,\bar y\mapsto\bar b]\) agree on
\(\{\bar x,\bar y\}\supseteq\FV(\varphi)\), so \((\ast)\) and
\cref{lem:fol-coincidence} give the claim.

\emph{Constants for parameters, forward direction.}  Let
\(X=\varphi(M^n,\bar b)\) with \(\bar b\in A^m\), and let
\(\psi\defeq\varphi[c_{b_1}/y_1]\cdots[c_{b_m}/y_m]\), the successive
substitutions of the constant terms \(c_{b_j}\) of
\cref{def:parameter-expansion}.  Since
\(\operatorname{Var}(c_{b_j})=\varnothing\), \cref{lem:substitution-fv}(1)
applied \(m\) times gives \(\FV(\psi)\subseteq\FV(\varphi)\setminus
\{\bar y\}\subseteq\{\bar x\}\), so \(\psi(\bar x;)\) is a declared
\(\Lang(A)\)-formula.  The value of \(c_{b_j}\) in \(\mathcal M_A\) is
\(b_j\) under every assignment, so \cref{thm:fol-substitution} applied \(m\)
times gives, for every assignment \(t\) in \(\mathcal M_A\),
\[
 \mathcal M_A,t\entails\psi\iff
 \mathcal M_A,t[y_m\mapsto b_m]\cdots[y_1\mapsto b_1]\entails\varphi
 \iff\mathcal M_A,t[\bar y\mapsto\bar b]\entails\varphi .
\]
By (F2), the right side is equivalent to
\(\mathcal M,t[\bar y\mapsto\bar b]\entails\varphi\), since
\(\mathcal M_A\mathbin{\upharpoonright}\Lang=\mathcal M\).  Taking
\(t=s[\bar x\mapsto\bar a]\) and using \((\ast)\) gives
\(\psi(M^n)=\varphi(M^n,\bar b)=X\), the left side computed in
\(\mathcal M_A\).

\emph{Constants for parameters, converse.}  Let \(\psi(\bar x;)\) be a
declared \(\Lang(A)\)-formula.  Only finitely many new constants occur in
\(\psi\); list them without repetition as \(c_{b_1},\ldots,c_{b_m}\), so
\(b_1,\ldots,b_m\in A\).  Choose distinct variables \(y_1,\ldots,y_m\)
outside \(\operatorname{Var}(\psi)\cup\{\bar x\}\).  For a term or formula
\(\theta\) of \(\Lang(A)\) built from symbols of \(\psi\), let
\(\theta^\circ\) be the result of replacing each occurrence of \(c_{b_j}\) by
\(y_j\): on terms by recursion, on atomic formulas termwise, commuting with
the connectives and with each quantifier, \((Qv\,\theta)^\circ=Qv\,
\theta^\circ\).  Then \(\varphi\defeq\psi^\circ\) is an \(\Lang\)-formula
with \(\FV(\varphi)\subseteq\FV(\psi)\cup\{\bar y\}\subseteq\{\bar x,\bar
y\}\).  We show, by induction on subformulas \(\theta\) of \(\psi\), that
for every assignment \(t\) in \(M\) with \(t(\bar y)=\bar b\),
\[
 \mathcal M_A,t\entails\theta\iff\mathcal M_A,t\entails\theta^\circ .
\]
For a term \(r\), the value of \(r^\circ\) under \(t\) equals that of \(r\),
by induction on \(r\), since \(t(y_j)=b_j=c_{b_j}^{\mathcal M_A}\).  This
settles atomic formulas, and the connective cases follow from the induction
hypothesis.  For \(\theta=Qv\,\theta_1\), the variable \(v\) occurs in
\(\psi\) and so is not any \(y_j\); hence each \(t[v\mapsto d]\) still sends
\(\bar y\) to \(\bar b\), and the induction hypothesis applies under every
such update.  Taking \(\theta=\psi\), \(t=s[\bar x\mapsto\bar a,\bar
y\mapsto\bar b]\), and using (F2) for the \(\Lang\)-formula \(\varphi\), we
get \(\psi(M^n)=\varphi(M^n,\bar b)\) with \(\varphi\) declared as
\((\bar x;\bar y)\).  So \(\psi(M^n)\) is \(A\)-definable in \(\mathcal M\).
\end{proof}
\end{annexitem}

\begin{annexitem}{annex:p03-w03-p03-003-closure}{thm:p03-w03-p03-003-closure}{Closure of the \(A\)-definable family}
\begin{proof}
Every verification below is carried out in an arbitrary
\(\Lang\)-structure \(\mathcal K\) and uses only the satisfaction clauses,
\cref{lem:fol-coincidence}, \((\ast)\) and (F1)--(F3).  Since each
constructed formula and tuple depends only on the given defining data and
\(\tau\), this proves the uniformity clause (6) together with (2)--(5).  When
\(\mathcal K=\mathcal M\) and the given tuples come from \(A\), the
constructed tuples also come from \(A\), as each is a concatenation of given
tuples and, in (5), of \(\bar c\in A^r\).

(1) Use the declared formulas \(x_1=x_1\) as \((x_1,\ldots,x_n;)\);
\(R(x_1,\ldots,x_p)\) as \((x_1,\ldots,x_p;)\);
\(f(x_1,\ldots,x_p)=x_{p+1}\) as \((x_1,\ldots,x_{p+1};)\), which for
\(p=0\) is \(c=x_1\) as \((x_1;)\); and \(x_i=x_j\) as
\((x_1,\ldots,x_n;)\).  The variables \(x_1,x_2,\ldots\) are distinct, each
list contains the free variables, and \((\ast)\) with the atomic
satisfaction clauses computes the solution sets as \(K^n\),
\(R^{\mathcal K}\), the graph of \(f^{\mathcal K}\), and the diagonal.

(2) The empty set is defined by \(\neg\,x_1=x_1\) as \((x_1,\ldots,x_n;)\).
If \(X=\varphi(K^n,\bar b)\) with \(\varphi(\bar x;\bar y)\), then
\(\neg\varphi\), declared with the same list, defines \(K^n\setminus X\) at
\(\bar b\), by \((\ast)\) and the clause for \(\neg\).  Let also
\(Y=\psi(K^n,\bar c)\) with \(\psi\) declared as \((\bar x';\bar y')\) and
\(\bar c\in K^{m'}\).  Choose \(m'\) distinct variables \(\bar z\) outside
\(\{\bar x,\bar y\}\) and, by the renaming clause of
\cref{thm:p03-w03-p03-003-normal-form}, a formula \(\psi'\) declared as
\((\bar x;\bar z)\) with \(\psi'(K^n,\bar c)=Y\).  Declare
\(\varphi\land\psi'\) and \(\varphi\lor\psi'\) as \((\bar x;\bar y\bar z)\),
a list of distinct variables containing both free-variable sets.  Under
\(s[\bar x\mapsto\bar a,\bar y\mapsto\bar b,\bar z\mapsto\bar c]\),
\cref{lem:fol-coincidence} evaluates \(\varphi\) as under
\(s[\bar x\mapsto\bar a,\bar y\mapsto\bar b]\) and \(\psi'\) as under
\(s[\bar x\mapsto\bar a,\bar z\mapsto\bar c]\); so the two formulas define
\(X\cap Y\) and \(X\cup Y\) at \(\bar b\bar c\).  Finite intersections and
unions follow by induction on the number of sets, the empty intersection
being \(K^n\) and the empty union \(\varnothing\).

(3) Let \(Y=\varphi(K^k,\bar b)\) with \(\varphi(x_1,\ldots,x_k;\bar y)\),
and choose distinct variables \(u_1,\ldots,u_n\) outside
\(\operatorname{Var}(\varphi)\cup\{\bar x,\bar y\}\).  Put
\[
 \theta\defeq\exists x_1\cdots\exists x_k
   \Bigl(\bigwedge_{i\leq k}x_i=u_{\tau(i)}\land\varphi\Bigr),
\]
declared as \((u_1,\ldots,u_n;\bar y)\); the list is distinct and
\(\FV(\theta)\subseteq(\{\bar x\}\cup\{\bar u\}\cup\{\bar y\})\setminus
\{\bar x\}=\{\bar u,\bar y\}\).  Let \(\bar a\in K^n\) and
\(t=s[\bar u\mapsto\bar a,\bar y\mapsto\bar b]\).  By (F1),
\(\mathcal K,t\entails\theta\) iff for some \(\bar e\in K^k\),
\(\mathcal K,t[\bar x\mapsto\bar e]\entails\bigwedge_ix_i=u_{\tau(i)}
\land\varphi\).  The update does not touch the \(u_l\), so by (F3) the
equations say \(e_i=a_{\tau(i)}\) for all \(i\), that is
\(\bar e=\tau_K(\bar a)\).  The update does not touch \(\bar y\) either, so
\(t[\bar x\mapsto\tau_K(\bar a)]\) sends \(\bar x\) to \(\tau_K(\bar a)\)
and \(\bar y\) to \(\bar b\), and by \((\ast)\) it satisfies \(\varphi\) iff
\(\tau_K(\bar a)\in Y\).  Hence
\(\theta(K^n,\bar b)=\tau_K^{-1}(Y)\).  The named special cases are
instances of \(\tau\), by \cref{def:p03-w03-p03-003-coordinate-maps}.

(4) Let \(X=\varphi(K^n,\bar b)\) with \(\varphi(x_1,\ldots,x_n;\bar y)\),
and choose distinct \(u_1,\ldots,u_k\) outside
\(\operatorname{Var}(\varphi)\cup\{\bar x,\bar y\}\).  Put
\[
 \theta\defeq\exists x_1\cdots\exists x_n
   \Bigl(\bigwedge_{i\leq k}u_i=x_{\tau(i)}\land\varphi\Bigr),
\]
declared as \((u_1,\ldots,u_k;\bar y)\); as in (3),
\(\FV(\theta)\subseteq\{\bar u,\bar y\}\).  For \(\bar d\in K^k\) and
\(t=s[\bar u\mapsto\bar d,\bar y\mapsto\bar b]\), (F1) and (F3) give:
\(\mathcal K,t\entails\theta\) iff there is \(\bar a\in K^n\) with
\(d_i=a_{\tau(i)}\) for every \(i\leq k\) and
\(\mathcal K,t[\bar x\mapsto\bar a]\entails\varphi\), that is, by \((\ast)\),
\(\bar a\in X\).  So \(\theta(K^k,\bar b)=\{\tau_K(\bar a):\bar a\in X\}
=\tau_K(X)\).  For projection onto the first \(k<n\) coordinates, the
formula \(\exists x_{k+1}\cdots\exists x_n\,\varphi\) declared as
\((x_1,\ldots,x_k;\bar y)\) also works, by (F1) and \((\ast)\) directly.

(5) Let \(X=\varphi(K^{n+r},\bar b)\) with
\(\varphi(x_1,\ldots,x_{n+r};\bar y)\), and let \(\bar c\in K^r\).  The list
\((x_1,\ldots,x_n;x_{n+1},\ldots,x_{n+r},\bar y)\) is the same set of
distinct variables, re-divided, so it is a declared list for \(\varphi\).
For \(\bar a\in K^n\) the assignment
\(s[(x_1,\ldots,x_n)\mapsto\bar a,(x_{n+1},\ldots,x_{n+r},\bar y)\mapsto
\bar c\bar b]\) equals
\(s[(x_1,\ldots,x_{n+r})\mapsto\bar a\bar c,\bar y\mapsto\bar b]\), so by
\((\ast)\) the new declared formula defines
\(\{\bar a:(\bar a,\bar c)\in X\}\) at \(\bar c\bar b\).
\end{proof}
\end{annexitem}

\begin{annexitem}{annex:p03-w03-p03-003-least}{thm:p03-w03-p03-003-least-family}{The \(A\)-definable sets form the least closed family}
\begin{proof}
\emph{Existence.}  The family \((\mathcal P(M^n))_{n\geq1}\) has
properties (1)--(5), and each property is preserved by componentwise
intersections of families having it, because each says that certain sets
belong to the family or that the family is closed under certain operations.
So the componentwise intersection \(\mathcal D\) of all such families has
the properties and is contained in each of them.

\emph{\(\mathcal D_n\subseteq\operatorname{Def}_n(\mathcal M/A)\).}  By
\cref{thm:p03-w03-p03-003-closure}, the family
\((\operatorname{Def}_n(\mathcal M/A))_{n\geq1}\) has properties (1)--(5),
so it contains the least such family.

\emph{\(\operatorname{Def}_n(\mathcal M/A)\subseteq\mathcal D_n\).}  For a
formula \(\theta\) and a list \(\bar v=(v_1,\ldots,v_k)\), \(k\geq1\), of
distinct variables with \(\FV(\theta)\subseteq\{\bar v\}\), put
\(\theta[\bar v]\defeq\{\bar c\in M^k:\mathcal M,s[\bar v\mapsto\bar c]
\entails\theta\}\), which is independent of \(s\) by
\cref{lem:fol-coincidence}.  For a term \(r\) with
\(\operatorname{Var}(r)\subseteq\{\bar v\}\), put
\[
 G_r[\bar v]\defeq\{(\bar c,d)\in M^{k+1}:
   \text{the value of }r\text{ under }s[\bar v\mapsto\bar c]\text{ is }d\},
\]
again independent of \(s\) by term coincidence.  We use that
\(\varnothing=M^k\setminus M^k\in\mathcal D_k\) and that the identity of
\(\{1,\ldots,k\}\) gives the identity coordinate map.

\emph{Terms.}  We show \(G_r[\bar v]\in\mathcal D_{k+1}\) by induction on
\(r\).  If \(r=v_i\), then \(G_r[\bar v]\) is the diagonal
\(\Delta^{k+1}_{i,k+1}\), which is in \(\mathcal D_{k+1}\) by (1).  Let
\(r=f(r_1,\ldots,r_p)\) with \(p\geq0\); each
\(\operatorname{Var}(r_j)\subseteq\{\bar v\}\), so
\(G_{r_j}[\bar v]\in\mathcal D_{k+1}\) by the induction hypothesis.  Index
the coordinates of \(M^{k+p+1}\) as \((\bar c,\bar e,d)\) with
\(\bar c\in M^k\), \(\bar e\in M^p\), \(d\in M\), and let \(\alpha_j\),
\(\beta\), \(\gamma\) be the coordinate maps
\((\bar c,\bar e,d)\mapsto(\bar c,e_j)\),
\((\bar c,\bar e,d)\mapsto(\bar e,d)\) and
\((\bar c,\bar e,d)\mapsto(\bar c,d)\).  Let \(\Gamma_f\subseteq M^{p+1}\)
be the graph of \(f^{\mathcal M}\), in \(\mathcal D_{p+1}\) by (1).  Then
\[
 G_r[\bar v]=\gamma\Bigl(\bigcap_{j\leq p}\alpha_j^{-1}(G_{r_j}[\bar v])
   \cap\beta^{-1}(\Gamma_f)\Bigr),
\]
where the intersection over \(j\) is \(M^{k+p+1}\) when \(p=0\).  Indeed
\((\bar c,d)\) lies in the right side iff there is \(\bar e\) such that each
\(e_j\) is the value of \(r_j\) under \(s[\bar v\mapsto\bar c]\) and
\(f^{\mathcal M}(\bar e)=d\), that is, iff \(d\) is the value of \(r\).  By
(2)--(4) the right side is in \(\mathcal D_{k+1}\).

\emph{Formulas.}  We show \(\theta[\bar v]\in\mathcal D_k\) for all
\(\theta\) and all admissible \(\bar v\), by induction on the complexity of
\(\theta\), for all lists simultaneously.  Throughout, write
\(\pi:M^{k+1}\to M^k\) for the coordinate map dropping the last coordinate.
\begin{itemize}
 \item \(\theta=\bottom\): \(\theta[\bar v]=\varnothing\).
 \item \(\theta=(r_1=r_2)\): \(\theta[\bar v]=\pi(G_{r_1}[\bar v]\cap
       G_{r_2}[\bar v])\), since \(\bar c\) is in it iff the two values agree,
       iff some \(d\) is the value of both.
 \item \(\theta=R(r_1,\ldots,r_p)\), \(p\geq1\): with coordinates
       \((\bar c,\bar e)\) on \(M^{k+p}\) and the coordinate maps
       \(\alpha_j:(\bar c,\bar e)\mapsto(\bar c,e_j)\),
       \(\beta:(\bar c,\bar e)\mapsto\bar e\),
       \(\pi':(\bar c,\bar e)\mapsto\bar c\),
       \[
        \theta[\bar v]=\pi'\Bigl(\bigcap_{j\leq p}
          \alpha_j^{-1}(G_{r_j}[\bar v])\cap\beta^{-1}(R^{\mathcal M})\Bigr),
       \]
       by the satisfaction clause for relation atoms.  If \(p=0\), then
       \(\theta[\bar v]\) is \(M^k\) when the empty tuple belongs to
       \(R^{\mathcal M}\) and \(\varnothing\) otherwise; both are in
       \(\mathcal D_k\).
 \item \(\theta=\neg\theta_1,\ \theta_1\land\theta_2,\ \theta_1\lor\theta_2,
       \ \theta_1\to\theta_2\): the components have free variables in
       \(\{\bar v\}\), and \(\theta[\bar v]\) is respectively
       \(M^k\setminus\theta_1[\bar v]\),
       \(\theta_1[\bar v]\cap\theta_2[\bar v]\),
       \(\theta_1[\bar v]\cup\theta_2[\bar v]\),
       \((M^k\setminus\theta_1[\bar v])\cup\theta_2[\bar v]\).
 \item \(\theta=\exists x\,\theta_1\).  If \(x\) is not among \(\bar v\),
       the list \(\bar vx\) is distinct and contains \(\FV(\theta_1)\), and
       \(s[\bar v\mapsto\bar c][x\mapsto d]=s[\bar vx\mapsto\bar cd]\); so
       \(\theta[\bar v]=\pi(\theta_1[\bar vx])\).  If \(x=v_i\), then
       \(\FV(\theta_1)\subseteq\{\bar v\}\) and
       \(s[\bar v\mapsto\bar c][v_i\mapsto d]=s[\bar v\mapsto\bar c']\), where
       \(\bar c'\) is \(\bar c\) with its \(i\)-th entry replaced by \(d\).
       Let \(\rho:\{1,\ldots,k\}\to\{1,\ldots,k+1\}\) be \(\rho(i)=k+1\) and
       \(\rho(j)=j\) for \(j\neq i\); then \(\rho_M(\bar c,d)=\bar c'\) and
       \(\theta[\bar v]=\pi(\rho_M^{-1}(\theta_1[\bar v]))\).
 \item \(\theta=\forall x\,\theta_1\).  With
       \(S=\theta_1[\bar vx]\) in the first case and
       \(S=\rho_M^{-1}(\theta_1[\bar v])\) in the second, the same
       computation with ``every'' gives
       \(\theta[\bar v]=M^k\setminus\pi(M^{k+1}\setminus S)\).
\end{itemize}
In every case the induction hypothesis puts the sets on the right in
\(\mathcal D\), and (2)--(4) put \(\theta[\bar v]\) in \(\mathcal D_k\).

\emph{Conclusion.}  Let \(X=\varphi(M^n,\bar b)\) with
\(\varphi(\bar x;\bar y)\) and \(\bar b\in A^m\).  If \(m=0\), then
\(X=\varphi[\bar x]\in\mathcal D_n\) by \((\ast)\).  If \(m\geq1\), then
\(\varphi[\bar x\bar y]\in\mathcal D_{n+m}\), and by \((\ast)\) \(X\) is its
fibre over \(\bar b\in A^m\), which lies in \(\mathcal D_n\) by (5).
\end{proof}
\end{annexitem}

\begin{annexitem}{annex:p03-w03-p03-003-auto}{thm:p03-w03-p03-003-automorphism}{Isomorphism transport and automorphism invariance}
\begin{proof}
The declared list \((\bar x;\bar y)\) is a list of \(n+m\) distinct
variables containing \(\FV(\varphi)\), so \cref{lem:iso-invariance} applies
to \(\varphi\) with these variables: for \(\bar a\in M^n\),
\(\mathcal M\entails\varphi(\bar a,\bar b)\) iff
\(\mathcal N\entails\varphi(h\bar a,h\bar b)\).  If
\(\bar a\in\varphi(M^n,\bar b)\), then \(h\bar a\in\varphi(N^n,h\bar b)\);
so \(h(\varphi(M^n,\bar b))\subseteq\varphi(N^n,h\bar b)\).  Conversely let
\(\bar a'\in\varphi(N^n,h\bar b)\).  As \(h\) is surjective,
\(\bar a'=h\bar a\) for some \(\bar a\in M^n\), and the equivalence gives
\(\bar a\in\varphi(M^n,\bar b)\).

For the second statement let \(X=\varphi(M^n,\bar b)\) with
\(\bar b\in A^m\) and \(\sigma\in\operatorname{Aut}(\mathcal M/A)\).  Then
\(\sigma\bar b=\bar b\), and the first statement with \(h=\sigma\) and
\(\mathcal N=\mathcal M\) gives \(\sigma(X)=\varphi(M^n,\sigma\bar b)=X\).
The contrapositive form is the same implication.
\end{proof}
\end{annexitem}

\begin{annexitem}{annex:p03-w03-p03-003-rigid}{thm:p03-w03-p03-003-rigid}{An invariant set that is not definable}
\begin{proof}
(1) Let \(\sigma\) be an automorphism of \((\N,<)\).  Its inverse is also an
automorphism, being a bijection that preserves and reflects \(<\).  We show
\(\sigma(n)=n\) by strong induction on \(n\).  Suppose \(\sigma(j)=j\) and
\(\sigma^{-1}(j)=j\) for all \(j<n\).  Since \(\sigma\) is injective and
\(\sigma(j)=j\) for \(j<n\), \(\sigma(n)\notin\{0,\ldots,n-1\}\), so
\(\sigma(n)\geq n\); likewise \(\sigma^{-1}(n)\geq n\).  If
\(\sigma(n)\neq n\), then \(\sigma^{-1}(n)\neq n\), so \(\sigma^{-1}(n)>n\),
and since \(\sigma\) preserves \(<\),
\(n=\sigma(\sigma^{-1}(n))>\sigma(n)\geq n\), a contradiction.  Hence
\(\sigma(n)=n\), and then also \(\sigma^{-1}(n)=n\).  So \(\sigma\) is the
identity.

(2) The symbols available for \(\Lang=\{<\}\) are \(<\), the logical and
punctuation tokens listed in \cref{sec:fol-syntax-semantics-01}, and the
variables \(x_0,x_1,\ldots\).  Assign the natural number \(i+12\) to
\(x_i\) and the numbers \(0,\ldots,11\) to the other twelve symbols.  Send a
finite sequence \((e_1,\ldots,e_\ell)\) of natural numbers to
\(2^{e_1+1}3^{e_2+1}\cdots p_\ell^{e_\ell+1}\), where \(p_\ell\) is the
\(\ell\)-th prime; by unique factorization this is injective on finite
sequences of every length.  A definable \(X\subseteq\N\) has defining data
consisting of a formula \(\varphi\), a declared list \((x;y_1,\ldots,y_m)\)
and a tuple \(\bar b\in\N^m\).  Code the data by the finite sequence
\((\ulcorner\varphi\urcorner,m,\#x,\#y_1,\ldots,\#y_m,b_1,\ldots,b_m)\),
where \(\ulcorner\varphi\urcorner\) codes the symbol string of \(\varphi\)
and \(\#v\) is the number of the variable \(v\).  The data can be read back
from the sequence, so the coding of data is injective.  Map each definable
\(X\subseteq\N\) to the least code of data defining it.  Different sets
have no common defining data, so this map is an injection into \(\N\), and
the set of definable subsets of \(\N\) is countable.  No choice principle is
used.

(3) For any \(F:\N\to\mathcal P(\N)\), the set
\(D=\{i\in\N:i\notin F(i)\}\) differs from every \(F(i)\) at \(i\), so
\(F\) is not surjective.  If every subset of \(\N\) were definable, the
injection \(g\) of (2) would be defined on all of \(\mathcal P(\N)\), and
\(F(i)=g^{-1}(i)\) for \(i\) in the range of \(g\), \(F(i)=\varnothing\)
otherwise, would be a surjection.  So some \(X\subseteq\N\) is not
definable.  By (1), \(\operatorname{Aut}(\mathcal M/A)\) contains only the
identity for every \(A\subseteq\N\), so \(\sigma(X)=X\) for each of its
members.
\end{proof}
\end{annexitem}

\begin{annexitem}{annex:p03-w03-p03-003-elem}{thm:p03-w03-p03-003-elementary-extension}{Transfer of definable sets to an elementary extension}
\begin{proof}
By the definition in \cref{sec:maps-elementarity-03}, for every formula
\(\chi\), every list \(\bar v\) of \(k\geq1\) distinct variables containing
\(\FV(\chi)\) and every \(\bar c\in M^k\),
\(\mathcal M,s[\bar v\mapsto\bar c]\entails\chi\) iff
\(\mathcal N,s'[\bar v\mapsto\bar c]\entails\chi\), for any assignments
\(s\) in \(M\) and \(s'\) in \(N\).

(1) For \(\bar a\in M^n\), apply this to \(\chi=\varphi\) and the list
\(\bar x\bar y\) with the tuple \(\bar a\bar b\): by \((\ast)\),
\(\bar a\in\varphi(N^n,\bar b)\) iff \(\bar a\in\varphi(M^n,\bar b)\).

(2) Let \(\varphi(\bar x;\bar y)\) and \(\psi(\bar x';\bar y')\), with
tuples \(\bar b\in M^m\) and \(\bar c\in M^{m'}\), define the same
\(X\subseteq M^n\).  Choose \(m'\) distinct variables \(\bar z\) outside
\(\{\bar x,\bar y\}\), and by the renaming clause of
\cref{thm:p03-w03-p03-003-normal-form} a formula \(\psi'\) declared as
\((\bar x;\bar z)\) with \(\psi'(K^n,\bar c)=\psi(K^n,\bar c)\) for both
\(\mathcal K=\mathcal M\) and \(\mathcal K=\mathcal N\).  Put
\[
 \chi\defeq\forall x_1\cdots\forall x_n\bigl((\varphi\to\psi')\land
   (\psi'\to\varphi)\bigr),
\]
so \(\FV(\chi)\subseteq\{\bar y,\bar z\}\).  For either structure
\(\mathcal K\) and \(t=s[\bar y\mapsto\bar b,\bar z\mapsto\bar c]\), (F1)
gives \(\mathcal K,t\entails\chi\) iff for every \(\bar a\in K^n\) the
assignment \(t[\bar x\mapsto\bar a]=s[\bar x\mapsto\bar a,\bar y\mapsto
\bar b,\bar z\mapsto\bar c]\) satisfies \(\varphi\) exactly when it
satisfies \(\psi'\); by \cref{lem:fol-coincidence} and \((\ast)\), iff
\(\varphi(K^n,\bar b)=\psi'(K^n,\bar c)\).  In \(\mathcal M\) the two sets
are both \(X\), so \(\mathcal M,t\entails\chi\).  If \(m+m'\geq1\),
elementarity for the list \(\bar y\bar z\) and the tuple \(\bar b\bar c\)
gives \(\mathcal N,s'[\bar y\mapsto\bar b,\bar z\mapsto\bar c]\entails\chi\).
If \(m=m'=0\), \(\chi\) is a sentence; choose any variable \(w\) and any
\(e\in M\) (domains are nonempty) and apply elementarity to the list
\((w)\) and the tuple \((e)\), using \cref{lem:fol-coincidence}.  In both
cases the computation above, made in \(\mathcal N\), gives
\(\varphi(N^n,\bar b)=\psi'(N^n,\bar c)=\psi(N^n,\bar c)\).

(3) By (2), \(X^{\mathcal N}\) does not depend on the defining data with
parameters in \(M\).  If \(X\) is \(A\)-definable, choosing data with
parameters in \(A\) shows \(X^{\mathcal N}\in\operatorname{Def}_n
(\mathcal N/A)\), and \(X^{\mathcal N}\cap M^n=X\) by (1).  The map is
injective: \(X^{\mathcal N}=Y^{\mathcal N}\) gives
\(X=X^{\mathcal N}\cap M^n=Y^{\mathcal N}\cap M^n=Y\).  It is surjective
onto \(\operatorname{Def}_n(\mathcal N/A)\): if \(Z=\varphi(N^n,\bar b)\)
with \(\bar b\in A^m\subseteq M^m\), then \(X=\varphi(M^n,\bar b)\) is
\(A\)-definable in \(\mathcal M\) and \(X^{\mathcal N}=Z\).  The formula
\(x_1=x_1\) declared as \((x_1,\ldots,x_n;)\) gives
\((M^n)^{\mathcal N}=N^n\).

For the operations, let the given sets be defined in \(\mathcal M\) by
some data with parameters in \(M\).  The same data define their extensions
in \(\mathcal N\), by the definition of \(X\mapsto X^{\mathcal N}\).  By the
uniformity clause of \cref{thm:p03-w03-p03-003-closure}, the formula and
tuple constructed there from these data define the result of the operation
in \(\mathcal M\), and the result of the same operation on the extensions in
\(\mathcal N\).  The first fact and (2) make the latter set the extension of
the result.  This gives
\((M^n\setminus X)^{\mathcal N}=N^n\setminus X^{\mathcal N}\),
\((X\cap Y)^{\mathcal N}=X^{\mathcal N}\cap Y^{\mathcal N}\),
\((X\cup Y)^{\mathcal N}=X^{\mathcal N}\cup Y^{\mathcal N}\), the two
identities for \(\tau\), and the statement on fibres; for the fibre over
\(\bar c\in A^r\) the constructed tuple is \(\bar c\) followed by the given
tuple, and \(\bar c\) is a tuple from \(M\).  Finally \(X\subseteq Y\) iff
\(X\cap Y=X\) iff \((X\cap Y)^{\mathcal N}=X^{\mathcal N}\), by
injectivity, iff \(X^{\mathcal N}\cap Y^{\mathcal N}=X^{\mathcal N}\) iff
\(X^{\mathcal N}\subseteq Y^{\mathcal N}\).

(4) The empty subset of \(M^n\) is defined by \(\neg\,x_1=x_1\) declared as
\((x_1,\ldots,x_n;)\), and the same formula defines the empty subset of
\(N^n\); so \(\varnothing^{\mathcal N}=\varnothing\), and by injectivity
\(X^{\mathcal N}=\varnothing\) only if \(X=\varnothing\).  Let
\(X=\{\bar a^1,\ldots,\bar a^r\}\) with \(r\geq1\).  Choose distinct
variables \(y^j_i\) (\(1\leq j\leq r\), \(1\leq i\leq n\)) outside
\(\{\bar x\}\), and declare
\[
 \varepsilon\defeq\bigvee_{j\leq r}\ \bigwedge_{i\leq n}x_i=y^j_i
\]
as \((\bar x;\bar y^1\cdots\bar y^r)\).  In any structure containing the
entries of the \(\bar a^j\), by (F3) and \((\ast)\),
\(\varepsilon\) at the parameter tuple \(\bar a^1\cdots\bar a^r\) defines
exactly \(\{\bar a^1,\ldots,\bar a^r\}\).  Applying this in \(\mathcal M\)
and in \(\mathcal N\) and using (2), \(X^{\mathcal N}=X\).  Conversely, if
\(X^{\mathcal N}\) is finite then so is \(X=X^{\mathcal N}\cap M^n\).
\end{proof}
\end{annexitem}
